%!TEX TS-program = xelatex
%!TEX encoding = UTF-8 Unicode
\documentclass[10pt,a4paper]{moderncv}

% 1. moderncv 主题
\moderncvtheme[blue]{classic}

% 2. 页面边距设置
\usepackage[scale=0.95]{geometry}
\setlength{\hintscolumnwidth}{3.0cm}
\recomputelengths % 更改页面布局后必须重新计算
\usepackage{enumitem}
\setlist[itemize]{leftmargin=*, itemsep=0.04em, topsep=0.08em, parsep=0em, partopsep=0em}

% 3. (核心) 设置字体
% 这一步是设置默认的 "英文字体"
% moderncv 默认会加载 lmodern (Latin Modern) 字体包，它们是高质量的英文字体
% 我们可以显式声明一下，或者不写这两行（让 moderncv 自动加载）
%\usepackage{tgheros}
%\setmainfont{TeX Gyre Heros}
%\setsansfont{TeX Gyre Heros}

% 4. (核心) 加载 xeCJK 来处理中文
\usepackage{xeCJK}
\XeTeXlinebreaklocale "zh" % 允许中文在正确的地方换行
\widowpenalty=10000
%\CJKtilde % 修正 CJK 字符和 ~ 之间的距离

% 5. (核心) 单独设置 "中文字体"
% 您之前设定的字体在这里使用
% 注意：请确保您的系统中安装了这两个字体，并且字体文件名无误
\IfFontExistsTF{Hiragino Sans GB}
  {\setCJKmainfont[Mapping=tex-text]{Hiragino Sans GB}\setCJKsansfont[Mapping=tex-text]{Hiragino Sans GB}}
  {\setCJKmainfont[Mapping=tex-text]{Noto Sans CJK SC}\setCJKsansfont[Mapping=tex-text]{Noto Sans CJK SC}}

% 6. 个人信息 (保持不变)
\firstname{}
\familyname{Buyu Deng}
\homepage{https://buyudeng.github.io/}
\email{buyudeng@stanford.edu}
\social[github]{BuyuDeng}
\address{Middle Campus, USTC}{230026 Hefei, Anhui}
\nopagenumbers{}
%\extrainfo{Born in December 2006}
\photo[60pt]{../assets/img/profile.jpg}                 % ‘64pt’是图片必须压缩至的高度、‘0.4pt‘是图片边框的宽度 (如不需要可调节至0pt)、’picture‘ 是图片文件的名字;可选项、如不需要可删除本行
%\quote{引言(可选项)}                           % 可选项、如不需要可删除本行

% 显示索引号;仅用于在简历中使用了引言
%\makeatletter
%\renewcommand*{\bibliographyitemlabel}{\@biblabel{\arabic{enumiv}}}
%\makeatother

% 分类索引
%\usepackage{multibib}
%\newcites{book,misc}{{Books},{Others}}
%----------------------------------------------------------------------------------
%            内容
%----------------------------------------------------------------------------------
\begin{document}
\maketitle
%----------------------------------------------------------------------------------
\section{\textbf{Education}}
\cventry
{2023.09 -- Expected 2027.06}
{B.S. Candidate}
{School of the Gifted Young}{}
{University of Science and Technology of China (USTC)}
{\textbf{Major:} Geophysics}
\vspace{0.45em}
\textbf{GPA:} 3.62/4.30 (\textbf{Rank}: 3/20)

\vspace{0.22em}
\textbf{Relevant Coursework:} Linear Algebra: \textbf{90}, Probability \& Statistics: \textbf{92}, Theoretical Mechanics: \textbf{95}, Python \& Deep Learning: \textbf{A}, Programming in Geophysics: \textbf{94}, Introduction to Seismic Exploration: \textbf{98}

\vspace{0.08em}
\textbf{Research Interests:} Seismic interpretation, geomodeling, forward modeling; deep learning, agentic AI
\vspace{0.16em}
%----------------------------------------------------------------------------------
\section{\textbf{Research Experience}}
\cventry{2026.07 -- Present}{Undergraduate Research Intern}{\href{https://scerf.stanford.edu/}{\textcolor{blue}{\underline{Stanford Center for Earth Resources Forecasting (SCERF)}}}}{}{Stanford University}{Supervisor: \textbf{Prof. Tapan Mukerji}}
\vspace{0.08em}
\begin{itemize}
    \item Conduct research on geomodeling and forward modeling for subsurface energy systems.
\end{itemize}
\vspace{0.12em}

\cventry{2024.08 -- Present}{Undergraduate Researcher}{\href{http://cig.ustc.edu.cn/}{\textcolor{blue}{\underline{Computational
Interpretation Group}}} (CIG)}{}{}
{Supervisor: \textbf{Prof. Xinming Wu}}
\vspace{0.45em}

\textbf{3D "All-in-One" Seismic Denoising}
\begin{itemize}
    \item Built a synthetic 3D seismic dataset with random noise, acquisition footprints, and migration artifacts; trained an all-in-one denoising network based on Degradation Representation Learning.\\
    Code: \href{https://github.com/BuyuDeng/NDR_Seismic}{\textcolor{blue}{\underline{github.com/BuyuDeng/NDR\_Seismic}}}
\end{itemize}
\vspace{0.12em}

% --- 第二个小标题：VAE压缩项目 ---
\textbf{Geo-VAE for 3D Geophysical Data Compression \& Latent Processing}
\begin{itemize}
    \item Adapted Wan-VAE with cached causal pseudo-temporal resampling along crossline, achieving isotropic 8$\times$8$\times$8 downsampling and 31.0:1 scalar compression for 256$^3$ cubes.
    \item Benchmarked five 3D geophysical modalities over 128$^3$--512$^3$ cubes; maintained competitive, modality-dependent reconstruction fidelity and reduced latency versus fine-tuned Wan-VAE at larger cube sizes.
    \item Developed RGT-conditioned seismic generation and 768$^3$ complete-volume latent denoising; one latent-map call achieved higher SSIM and lower leakage than raw-space controls.\\
    Full manuscript under revision at \textit{Geophysics}; gave an oral presentation at IMAGE 2026.\\
    Code: \href{https://github.com/BuyuDeng/Geo-VAE}{\textcolor{blue}{\underline{github.com/BuyuDeng/Geo-VAE}}}
\end{itemize}
\vspace{0.12em}

\textbf{PHY-to-BGC Ocean Biogeochemical Emulation}
\begin{itemize}
    \item Built a 2.5D full-image U-Net emulator mapping global physical ocean fields to coupled biogeochemical variables, with evaluation under anomalous ocean-climate conditions; \textbf{manuscript in preparation}.
\end{itemize}
\vspace{0.12em}

\textbf{GNSS-NISAR Ionospheric Super-Resolution Reconstruction}
\begin{itemize}
    \item Developed an automated pipeline aligning GNSS TEC-derived ionospheric phase with NISAR dual-frequency InSAR for high-resolution super-resolution research.
\end{itemize}
%----------------------------------------------------------------------------------

\section{\textbf{Honors and Awards}}

\cvitem{2025}{\textbf{Zhao Jiuzhang Elite Program Scholarship}, School of Earth and Space Sciences}
\cvitem{2024}{\textbf{Rose Fund Endeavor Scholarship}, School of the Gifted Young}
\cvitem{2024}{\textbf{Zhao Jiuzhang Elite Program Scholarship}, School of Earth and Space Sciences}
\cvitem{2023}{\textbf{Freshman Scholarship}, University of Science and Technology of China (USTC)}
%----------------------------------------------------------------------------------

\section{\textbf{Experience and Leadership}}
\cvitem{2025.08}{\textbf{IPGP Geophysical Field Camp}: conducted geophysical surveys and joined academic seminars with faculty from the Institut de Physique du Globe de Paris (IPGP).}
\cvitem{2025.06 -- 2026.01}{\textbf{Head (Minister)}, Dept. of Cultural \& Sports in College Student Union}
\end{document}
